\documentclass[conference]{IEEEtran}

\usepackage[utf8]{inputenc}
\usepackage[T1]{fontenc}
\usepackage[spanish]{babel}
\usepackage{graphicx}
\usepackage{booktabs}
\usepackage{url}
\usepackage{amsmath}

\newcommand{\placeholderfigure}[2]{%
    \IfFileExists{#1}{%
        \includegraphics[width=0.95\linewidth]{#1}%
    }{%
        \fbox{\parbox{0.90\linewidth}{\centering #2}}%
    }%
}

\title{Sistema de Mapeo, Navegacion Autonoma y Busqueda de Objetivos en TurtleBot}

\author{
\IEEEauthorblockN{TODO: integrantes del grupo}
\IEEEauthorblockA{I402 - Principios de la Robotica Autonoma\\
TODO: institucion y correos}
}

\begin{document}
\maketitle

\begin{abstract}
Se presenta una primera version del sistema desarrollado para el trabajo final de
robotica autonoma, orientado a construir un mapa de ocupacion, navegar sobre el
mismo y desplegar una mision de busqueda de conos rojos en un TurtleBot real. La
solucion integra un modulo de simultaneous localization and mapping (SLAM)
basado en grillas y filtros de particulas, una pila de navegacion con Monte Carlo
localization (MCL), planificacion A*, seguimiento cinematico y una finite state
machine (FSM) de mision. Para el despliegue en hardware se incorporaron perfiles
de TurtleBot4, remapeos de topicos, ajustes de transformadas, deteccion visual de
conos y un kit de evidencia reproducible con rosbags, registros estructurados,
replay en RViz y reportes post-ejecucion. El estado actual incluye validaciones
por codigo, pruebas unitarias y smokes sinteticos; la validacion end-to-end sobre
el robot real queda como TODO experimental para completar con la corrida de
laboratorio.
\end{abstract}

\begin{IEEEkeywords}
robotica movil, slam, navegacion autonoma, turtlebot, sim-to-real
\end{IEEEkeywords}

\section{Introduccion}

El trabajo final propone integrar percepcion, localizacion, planificacion y
control en un sistema autonomo para un robot movil. La consigna se organiza en
tres partes encadenadas: la Parte A construye un mapa del entorno, la Parte B usa
ese mapa para navegacion autonoma y la Parte C traslada el sistema al robot real
para buscar y alcanzar conos rojos en un laberinto fisico.

En este informe se documenta una arquitectura basada en mapa grillado, sin
landmarks como insumo central. Esta eleccion corresponde al Resultado V1 de la
consigna: un mapa de ocupacion que alimenta una navegacion por grilla. La
Fig.~\ref{fig:arquitectura-general} resume el flujo implementado.

\begin{figure}[t]
    \centering
    % TODO: reemplazar por diagrama exportado a docs/informe/figures/fig_arquitectura_general.png
    \placeholderfigure{figures/fig_arquitectura_general.png}{Placeholder: arquitectura general, sensores, SLAM, MCL, A*, percepcion, FSM y evidencia.}
    \caption{Arquitectura general del sistema, desde sensores y mapa hasta navegacion, percepcion de cono y evidencia reproducible.}
    \label{fig:arquitectura-general}
\end{figure}

\section{Alcance y Criterios de Consigna}

La Parte A exige implementar SLAM y obtener un mapa usable para navegacion. La
Parte B exige localizacion probabilistica, planificacion libre de colisiones,
seguimiento de trayectoria, orientacion final, replanning y respuesta ante
obstaculos no mapeados. La Parte C exige ejecutar sobre robot real, detectar
conos rojos, ignorar distractores, planificar caminos validos hacia el objetivo y
analizar la brecha sim-to-real.

Un punto critico de la consigna es que una deteccion visual no debe interpretarse
como camino directo libre. Si el cono se observa a traves de una abertura o
rejilla, el sistema debe validar el objetivo contra el mapa y no atravesar
paredes. Esta restriccion se incorporo en la FSM de mision mediante validacion de
goals antes de publicar comandos de navegacion.

\section{Arquitectura Implementada}

La implementacion se divide en paquetes ROS 2. El paquete \texttt{maze\_slam}
construye el mapa; \texttt{maze\_nav} publica el mapa, localiza al robot y navega;
\texttt{maze\_perception} detecta conos en la imagen; y \texttt{maze\_mission}
coordina la busqueda y aproximacion al objetivo. La Tabla~\ref{tab:modulos}
resume las responsabilidades principales.

\begin{table}[t]
\centering
\caption{Modulos principales del sistema.}
\label{tab:modulos}
\begin{tabular}{@{}lll@{}}
\toprule
Modulo & Archivo principal & Responsabilidad \\
\midrule
SLAM & \texttt{fastslam\_node.py} & mapa y pose corregida \\
MCL & \texttt{localizer.py} & pose sobre mapa \\
Nav & \texttt{navigator.py} & A*, path following, recovery \\
Vision & \texttt{cone\_detector\_node.py} & deteccion de cono rojo \\
Mision & \texttt{mission\_node.py} & FSM y goals validados \\
\bottomrule
\end{tabular}
\end{table}

\section{Parte A: SLAM y Mapa}

Se implemento un Grid-Based FastSLAM en \texttt{maze\_slam}. Cada particula
mantiene su propia pose y un mapa de ocupacion en log-odds. Para cada scan se
aplica prediccion por odometria, ajuste local de pose, ponderacion por likelihood
field, integracion del LIDAR y remuestreo cuando corresponde.

El nodo \texttt{fastslam\_node.py} consume LIDAR y odometria, publica
\texttt{/map}, \texttt{/belief}, \texttt{/maze\_slam/particles} y la transformada
\texttt{map->odom}. El mapa real disponible para las pruebas de laboratorio es
\texttt{maps/laberinto\_lab\_20260702.yaml}, con resolucion de 0.03~m/celda,
origen \((-7.5,-7.5)\) y grilla de \(500 \times 500\) celdas. La
Fig.~\ref{fig:mapa-generado} debera reemplazarse por la version final usada en
la defensa.

\begin{figure}[t]
    \centering
    % TODO: insertar mapa exportado/overlay desde maps/laberinto_lab_20260702.pgm
    \placeholderfigure{figures/fig_mapa_laberinto.png}{Placeholder: mapa grillado del laberinto real usado por Parte B y Parte C.}
    \caption{Mapa de ocupacion del laberinto real utilizado como insumo para localizacion y planificacion.}
    \label{fig:mapa-generado}
\end{figure}

\section{Parte B: Navegacion Autonoma}

La navegacion se implemento en \texttt{maze\_nav}. La localizacion usa Monte
Carlo localization (MCL), con particulas inicializadas desde
\texttt{/initialpose}, actualizacion por LIDAR y publicacion de
\texttt{/amcl\_pose} y \texttt{/particlecloud}. En el robot real, el launch
\texttt{nav\_tb4\_live.launch.py} remapea \texttt{scan\_topic} y
\texttt{odom\_topic} hacia \texttt{/tb4\_0/scan} y \texttt{/tb4\_0/odom}, y
ajusta el montaje del LIDAR mediante \texttt{scan\_yaw\_offset} y
\texttt{scan\_x\_offset}.

El planificador global usa A* sobre una grilla inflada. El costo incluye una
penalizacion por proximidad a obstaculos para favorecer caminos centrados en los
pasillos. El seguimiento del camino usa un punto de lookahead y comandos
\texttt{/cmd\_vel}. El estado de navegacion se publica en \texttt{/nav\_state}.
La Fig.~\ref{fig:flujo-navegacion} muestra el flujo previsto de navegacion.

\begin{figure}[t]
    \centering
    % TODO: reemplazar por diagrama de mapa, costmap, A* y follower
    \placeholderfigure{figures/fig_flujo_navegacion.png}{Placeholder: mapa -> costmap -> A* -> path -> follower -> cmd_vel.}
    \caption{Flujo de navegacion: localizacion, costmap inflado, planificacion A* y seguimiento de trayectoria.}
    \label{fig:flujo-navegacion}
\end{figure}

Para el laberinto real se observo que la inflacion total de 0.26~m cerraba
pasillos angostos. Por ese motivo se ajusto la combinacion
\texttt{robot\_radius=0.14} e \texttt{inflation=0.06}, equivalente a 0.20~m. Este
valor debe coincidir con la validacion de goals en la mision y con la generacion
de waypoints.

\begin{figure}[t]
    \centering
    % TODO: captura de RViz con /map, /particlecloud, /plan y goal real
    \placeholderfigure{figures/fig_rviz_goal_nav.png}{Placeholder: RViz con mapa, particulas MCL, goal y path planificado.}
    \caption{Ejemplo de navegacion hacia un objetivo en RViz, mostrando mapa, particulas, goal y camino planificado.}
    \label{fig:rviz-goal-nav}
\end{figure}

\subsection{Obstaculos No Mapeados}

La consigna exige manejar obstaculos que no estaban presentes en el mapa. Para
esto el navegador calcula el clearance frontal y, si cae por debajo de
\texttt{safety\_stop}, entra en \texttt{RECOVERY}. En ese estado frena, ejecuta
un backoff corto, intenta registrar obstaculos dinamicos y replantea el camino.
Si el bloqueo persiste, publica \texttt{cmd\_vel=0} y reporta la razon en
\texttt{/nav\_debug}. La Fig.~\ref{fig:obstaculo-nav-debug} debe completarse con
una captura de replay o de laboratorio.

\begin{figure}[t]
    \centering
    % TODO: reemplazar por replay con obstaculo tipo silla/patas y nav_debug visible
    \placeholderfigure{figures/fig_obstaculo_nav_debug.png}{Placeholder: replay con obstaculo no mapeado, RECOVERY/BLOCKED y razon de nav_debug.}
    \caption{Diagnostico de obstaculo no mapeado mediante replay RViz y telemetria \texttt{/nav\_debug}.}
    \label{fig:obstaculo-nav-debug}
\end{figure}

\section{Parte C: Robot Real y Cono Rojo}

La Parte C agrega percepcion visual y una FSM de mision. El detector de cono
segmenta rojo en HSV, extrae blobs y publica detecciones estructuradas con
bearing, area y confianza. La posicion metrica del cono se estima combinando el
bearing de camara con el rango del LIDAR.

La FSM se implemento en \texttt{mission\_node.py}. Sus estados principales son
\texttt{LOCALIZE}, \texttt{SEARCH\_CONE}, \texttt{CONE\_DETECTED},
\texttt{ESTIMATE\_CONE\_GOAL}, \texttt{PLAN\_TO\_CONE},
\texttt{NAVIGATE\_TO\_CONE}, \texttt{AVOID\_OBSTACLE}, \texttt{REPLAN},
\texttt{VERIFY\_CONE}, \texttt{DONE} y \texttt{FAILURE}. La FSM no controla
velocidades: emite goals por \texttt{/goal\_pose} y observa el estado del
navegador.

\begin{figure}[t]
    \centering
    % TODO: diagrama de estados exportado a PNG
    \placeholderfigure{figures/fig_fsm_mision_cono.png}{Placeholder: FSM de mision del cono, desde LOCALIZE hasta DONE/FAILURE.}
    \caption{Maquina de estados de la mision de busqueda y aproximacion al cono rojo.}
    \label{fig:fsm-cono}
\end{figure}

La deteccion visual del cono no se usa como prueba de camino libre. Antes de
publicar un objetivo, \texttt{mission\_node.py} valida que el punto estimado no
caiga sobre una pared del mapa crudo y que el goal final sea navegable en el mapa
inflado. Este mecanismo cubre el caso en que el cono se observe a traves de una
rejilla o abertura. La Fig.~\ref{fig:detector-cono} debe completarse con una
captura del detector en el laboratorio.

\begin{figure}[t]
    \centering
    % TODO: captura real de /cone_debug_image y /cone_mask
    \placeholderfigure{figures/fig_detector_cono.png}{Placeholder: imagen anotada del detector de cono y mascara HSV.}
    \caption{Salida de depuracion del detector de cono rojo, con bounding box, centroide, bearing y mascara.}
    \label{fig:detector-cono}
\end{figure}

\section{Evidencia Experimental y Reproducibilidad}

Para evitar depender solo de video de celular, se incorporo un kit de evidencia
que registra rosbags, eventos estructurados, poses, goals, estados, velocidades,
detecciones, ultimo path y reportes. El script principal es
\texttt{scripts/lab\_record\_all.sh}; el postprocesamiento se realiza con
\texttt{scripts/lab\_make\_report.py}. El replay con visualizacion se apoya en
\texttt{scripts/lab\_replay\_rviz.sh}, \texttt{scripts/lab\_viz\_markers.py} y
\texttt{rviz/lab\_replay.rviz}.

\begin{figure}[t]
    \centering
    % TODO: reemplazar por captura real generada desde lab_replay_rviz.sh
    \placeholderfigure{figures/fig_rviz_replay_real.png}{Placeholder: replay RViz con mapa, trayectoria, path, goal, NAV/MISSION y deteccion.}
    \caption{Reconstruccion de una corrida real mediante replay de RViz, mostrando mapa, trayectoria estimada, objetivo, camino planificado y estados.}
    \label{fig:rviz-replay}
\end{figure}

\section{Resultados Preliminares}

Hasta el momento se cuenta con evidencia de codigo, documentacion y smokes. Se
confirmo que el mapa real esta versionado, que los launchs reales aceptan
namespace de TurtleBot4, que la navegacion remapea \texttt{/cmd\_vel} al topico
del robot y que el localizador detecta el frame de odometria real. Tambien se
documentaron smokes de la FSM con cono mockeado y un smoke de obstaculo frontal
persistente.

La Tabla~\ref{tab:estado-validacion} separa lo validado de lo pendiente. No se
afirma aun que el robot haya completado la mision real del cono.

\begin{table}[t]
\centering
\caption{Estado preliminar de validacion.}
\label{tab:estado-validacion}
\begin{tabular}{@{}ll@{}}
\toprule
Item & Estado \\
\midrule
Mapa real versionado & confirmado \\
Build temporal limpio & documentado \\
Tests maze\_mission/perception & documentado \\
FSM reachable/wall & smoke documentado \\
Obstaculo tipo silla & smoke sintetico \\
Mision real completa & TODO laboratorio \\
Silla real & TODO laboratorio \\
Convergencia inicial sin teleop & TODO laboratorio \\
\bottomrule
\end{tabular}
\end{table}

\section{Analisis Sim-to-Real}

El despliegue en hardware introdujo diferencias respecto de simulacion. Entre
ellas se observaron topicos namespaced, necesidad de publicar en
\texttt{/tb4\_0/cmd\_vel}, diferencias de TF, montaje fisico del LIDAR, pasillos
mas angostos que los margenes iniciales y dependencia de movimiento para que la
MCL actualice. Estas diferencias se trataron con remapeos, parametros de sensor,
reduccion conservadora de inflacion y documentacion de pendientes. Si durante la
corrida final el sistema falla parcialmente, se reportara la causa con rosbag,
\texttt{/nav\_debug}, timeline y replay.

\section{Conclusiones Preliminares}

Se integro una solucion modular que cubre las responsabilidades principales de
la consigna: mapeo, localizacion probabilistica, planificacion, seguimiento,
respuesta ante obstaculos, percepcion visual y supervision de mision. El estado
actual es defendible como primera version tecnica, pero requiere completar la
evidencia de laboratorio: corrida real, replay RViz, detector de cono bajo
iluminacion real, caso de obstaculo no mapeado y analisis de los errores
observados.

\begin{thebibliography}{9}

\bibitem{thrun2005}
S. Thrun, W. Burgard, and D. Fox, \emph{Probabilistic Robotics}. MIT Press, 2005.

\bibitem{ros2}
Open Robotics, ``ROS 2 Documentation,'' \url{https://docs.ros.org/}.

\bibitem{turtlebot4}
Clearpath Robotics, ``TurtleBot 4 User Manual,'' \url{https://turtlebot.github.io/turtlebot4-user-manual/}.

\end{thebibliography}

\end{document}

